\section{Ambrose and Augustine in Milan}

Between the autumn of 384 and Easter of 387, Milan brought together two of the four Latin Doctors of the Church. Viewed from Augustine's life, those thirty months were the decisive pivot from Manichaean disillusionment to Catholic baptism. Viewed from Ambrose's evidentiary horizon, the encounter looks quieter and more asymmetric---and for that very reason, deeply illuminating about how Ambrose actually pastored his city.

In the entire surviving corpus of Ambrose's letters and treatises, Augustine never appears by name. In 384, Ambrose was a forty-four-year-old metropolitan fighting the Altar of Victory debate, negotiating with a usurper, and bracing for an imperial assault on his churches; Augustine was a newly appointed thirty-year-old municipal professor of rhetoric, sent up from Rome by Symmachus himself (\work{Confessions} V.13.23). Yet the portrait Augustine later drew of Ambrose in \work{Confessions} V--IX and in his \work{Letters} remains the sharpest personal eyewitness sketch of Ambrose we possess.

Ambrose greeted the newcomer ``like a father, and looked with a benevolent and episcopal kindliness on my change of abode'' (\work{Confessions} V.13.23). Augustine began attending Sunday sermons as a connoisseur of style, judging whether the bishop's fame matched his delivery, and finding Ambrose's speech ``more erudite, yet less cheerful and soothing in manner, than that of Faustus'' the Manichaean. Soon, however, the substance entered with the rhetoric. What dismantled Augustine's Manichaean caricature of the Old Testament was Ambrose's repeated homiletic use of 2~Corinthians 3:6:

\begin{quote}
And with delight I heard Ambrose, in his sermons to the people, oftentimes most diligently recommend this text as a rule,---``The letter killeth, but the Spirit giveth life;'' whilst, drawing aside the mystic veil, he spiritually laid open that which, accepted according to the ``letter,'' seemed to teach perverse doctrines. (\work{Confessions} VI.4.6)
\end{quote}

Through those Sunday sermons, Augustine also learned that Catholic teaching on man created in the image of God did not conceive God as bounded by a human physical body (\work{Confessions} VI.3.4).

Just as vivid is Augustine's glimpse of Ambrose's daily habits. Ambrose kept an open door; visitors entered unannounced, yet crowds of troubled citizens left him almost no private leisure. In the rare minutes between pastoral business and necessary meals, Augustine watched him studying silently:

\begin{quote}
But while reading, his eyes glanced over the pages, and his heart searched out the sense, but his voice and tongue were silent. Ofttimes, when we had come (for no one was forbidden to enter, nor was it his custom that the arrival of those who came should be announced to him), we saw him thus reading to himself, and never otherwise; and, having long sat in silence (for who durst interrupt one so intent?), we were fain to depart. (\work{Confessions} VI.3.3)
\end{quote}

In an ancient world where reading aloud was ordinary practice, Ambrose's silent reading struck Augustine as remarkable; Augustine wondered whether the bishop read silently to spare a fragile voice or to avoid being drawn into impromptu textual debates during his brief hours of study. Because Ambrose was so absorbed in civic and diocesan burdens, Augustine never secured a long private seminar with him; instead, he turned for unhurried spiritual counsel to Ambrose's own mentor, the elderly presbyter Simplicianus (\work{Confessions} VIII.2.3).

Ambrose's pastoral relationship with Augustine's mother Monica was warmer and more direct. When Monica arrived in Milan carrying the North African custom of bringing cakes, bread, and diluted wine to the shrines of the martyrs, the doorkeeper turned her away on Ambrose's orders. Ambrose had prohibited such martyr-shrine banquets across Milan lest they foster drunkenness and resemble pagan \latin{parentalia} for the dead. Monica submitted immediately out of devotion to Ambrose, and Ambrose frequently praised her fervent piety to Augustine, congratulating him on having such a mother (\work{Confessions} VI.2.2).

A second exchange over Monica's scruples gave Christian history one of its most famous practical maxims. Perplexed because the Church of Milan did not fast on Saturdays whereas Rome and her own African town did, Monica asked Augustine to consult Ambrose. In two separate letters (\work{Letter} 36.14.32 to Casulanus and \work{Letter} 54.2.3 to Januarius), Augustine preserved Ambrose's answer verbatim---calling it an ``oracle from heaven'':

\begin{quote}
He answered, ``What else can I recommend to others than what I do myself?'' When I thought that by this he intended simply to prescribe to us that we should take food on Saturdays---for I knew this to be his own practice---he, following me, added these words: ``When I am here I do not fast on Saturday; but when I am at Rome I do: whatever church you may come to, conform to its custom, if you would avoid either receiving or giving offence.'' (St.~Augustine, \work{Letter} 36.14.32)
\end{quote}

After his conversion in late summer 386 and retreat to Cassiciacum, Augustine wrote to Ambrose notifying him of his past errors and asking which book of Scripture to read before baptism. Ambrose recommended the prophet Isaiah---an inspired pastoral choice for its prophetic breadth, though Augustine confessed that he found the opening chapters impenetrable and set the scroll aside for later (\work{Confessions} IX.5.13). In Lent 387, Augustine, his friend Alypius, and his teenage son Adeodatus returned to Milan, enrolled for baptism, and received the sacrament from Ambrose at Easter 387 (\work{Confessions} IX.6.14; by historical Easter computation, late April 387, with Easter Sunday falling on 25~April).

Augustine's debt to Ambrose lasted for the rest of his life. In \work{On Christian Doctrine} IV.21.46--50, he chose Ambrose alongside Cyprian to illustrate all three registers of Christian eloquence---quoting \work{De Spiritu Sancto} for the subdued style and \work{De Virginibus} for the temperate and grand styles. Years later, during the Pelagian controversy, when Pelagius tried to enlist Ambrose on his own side---calling him ``the blessed Bishop Ambrose, in whose writings the Roman faith shines forth with especial brightness, and whom the Latins have always regarded as the very flower and glory of their authors''---Augustine answered by quoting Ambrose's \work{Exposition on Luke} right back at Pelagius: ``no man is able to build without the Lord, no man to watch without the Lord, no man to undertake anything without the Lord'' (\work{On the Grace of Christ} I.43--44 [ch.~47--48]). In Milan, Ambrose had baptized the student who would become the foremost interpreter of his teaching.
